\documentclass[11pt]{article}
\usepackage[T1]{fontenc}
\usepackage[margin=0.95in]{geometry}
\usepackage{amsmath,amssymb,amsthm,booktabs,array,hyperref,needspace}
\hypersetup{colorlinks=true,linkcolor=blue,citecolor=blue,urlcolor=blue}
\emergencystretch=1.5em
\newtheorem{theorem}{Theorem}
\newtheorem{lemma}[theorem]{Lemma}
\newtheorem{proposition}[theorem]{Proposition}
\newtheorem{corollary}[theorem]{Corollary}
\newcommand{\Tr}{\operatorname{Tr}}
\newcommand{\Rcal}{\mathcal R_{\lambda,\tau}}
\newcommand{\Ecal}{\mathcal E_\lambda}
\newcommand{\id}{\operatorname{id}}
\title{Constant-gap entanglement distillation\\and a spectral criterion for noise-universal communication}
\author{Lluis Eriksson}
\date{18 October 2026 --- Working paper}
\begin{document}\maketitle
\begin{abstract}
We prove dimension-independent additive bounds on two-way entanglement
distillation. For every conjugate local diagonal unitary invariant
state, distillable entanglement, secret key and relative entropy of
entanglement lie in an explicitly computable interval of width at most
$\log_2 5$ bits per copy. For the broader class invariant under local
diagonal signs, the diagonal-relative-entropy reference exceeds
distillable entanglement by at most $\log_2 9$.
The proof uses established partial breeding with a stopped classification
whose entropy is bounded independently of dimension; a reduction-based
seed pays for all initial assistance. The resulting recovery bound
also implies at most $\log_2 6$ ebits of irreversibility for any mixture
of a pure state and separable noise, without a marginal constraint,
and controls mixtures with a mixed maximally correlated component.
For an identity component mixed with arbitrary channel noise, we give
six capacity guarantees and reduce their simultaneous worst cases to
depolarizing channels. An explicit isotropic interval has absolute width
tending to $2(1-f)$ at fixed entangled fraction $f$, refining the known
relative large-dimension approximation. For general discrete Hamiltonians with a finite partition sum at some
temperature, we identify the exact spectral boundary of absolute
entanglement-assisted recovery: for each fixed reset state, the loss
vanishes precisely when the partition sum diverges at its critical
inverse temperature. In that regime all six capacities attain their
erasure benchmarks with absolute $o(1)$ loss. In the complementary
regime every fixed reset has a strictly positive limiting assisted
deficit, given by a critical-Gibbs variational formula. An explicit
integer spectrum has an exact constant-loss plateau. The reset may
have infinite energy or entropy. Oscillator consequences include
reversible distributed states and a vanishing classical arbitrary-partner
additivity defect. Optimal constants, efficient implementations and
general exact finite-energy capacities remain unresolved.
\end{abstract}

\section{Introduction and scope}
We give an explicit, dimension-independent additive approximation to
entanglement distillation for two established symmetry classes of
bipartite quantum states. The computable lower endpoint uses the
coherent block on $\operatorname{span}\{|jj\rangle\}$; a stopped
classification isolates that block at bounded entanglement cost.
For conjugate local diagonal unitary invariant (CLDUI) states, the
distillable entanglement, secret key and relative entropy of entanglement
lie in one interval of width at most $\log_2 5$ bits per copy.
For local diagonal orthogonal invariant (LDOI) states, an explicit
diagonal-relative-entropy reference is within $\log_2 9$ of
distillable entanglement. These are additive bounds, rather than
approximations normalized by $\log d$.

The invariant classes and their block characterizations are established
\cite{jm,sn}. Recent semidefinite hierarchies study their separability
structure \cite{bl}; our claim concerns operational distillation bounds,
not the introduction or separability characterization of those classes.
The underlying parity measurements, partial breeding, adaptive
discarding and removal of assistance also have established antecedents
\cite{hostens,vv,maneva,leung}. Vollbrecht--Wolf
\cite[Section VI.B--C]{vw} obtained optimal large-dimension rates for
isotropic states after normalization by $\log d$. We use those protocol
primitives, hashing \cite{dw}, reduction-based distillability \cite{hh}
and the logarithmic reduction argument \cite{vwentropy} with explicit
attribution. The contribution is the bounded-entropy stopped
classification and its absolute-loss consequences, not those primitives.

The same argument proves that mixing any pure state with arbitrary
separable noise creates at most $\log_2 6$ ebits of irreversibility,
with no matching-marginal assumption. It also treats a mixed maximally
correlated useful component, for which the controlled quantities are
$E_D,K_D,E_R^\infty,E_R$, rather than entanglement cost.
Corollary~\ref{cor:ebchannel} applies the pure-state result to every
identity/entanglement-breaking channel mixture. Section~\ref{sec:minimax}
then allows arbitrary input-dependent noise: standard twirling organizes
six simultaneous worst-case capacities, while the classification theorem
provides a new absolute two-way interval for the isotropic minimizer.
The minimax reduction does not solve its unknown capacity functions.
The final part applies the recovery bound to oscillator reset channels,
obtaining absolute vanishing gaps for six energy-constrained capacities
for any fixed reset state, including infinite energy or entropy.
We give the finite-dimensional operational proof first, then the
oscillator construction and its separate channel-theoretic dependencies.

Section~\ref{sec:spectral} determines when the oscillator mechanism
extends to a general discrete Hamiltonian. The critical Gibbs entropy
description is established \cite[Proposition 1a]{shirokoventropy};
we derive an operational dichotomy from it. Divergence of the partition
sum at the critical inverse temperature permits entropy-preserving
escape and six absolute recovery laws. If the critical sum is finite,
every fixed reset state instead has a positive entanglement-assisted
deficit, characterized by relative entropy to the critical Gibbs state
plus a continuous binary-flag loss. The proof uses infinite-dimensional
relative-entropy sufficiency \cite{jencova} and finite-flag continuity
\cite{shirokovcontinuity}. A constructed spectrum makes the loss
constant at every sufficiently large energy. The scalar vacuum-reset
decomposition is credited to \cite[Appendix C]{catalano}; the new claim
is the capacity criterion and defect, not that decomposition or the
critical entropy description.

The constants are not claimed optimal, and the protocols are asymptotic
existence results, not efficient coding algorithms. We do not claim
general mixed-state reversibility, resolution of NPT bound entanglement,
general exact finite-energy reset capacities, or exhaustive priority certification.

\section{Entanglement recovery with a constant absolute loss}\label{sec:recovery}
All systems in this section are finite dimensional. $E_D$ is two-way
distillable entanglement, $E_C$ is standard LOCC entanglement cost, and
$K_D$ is distillable secret key, with no initial entanglement or key and
free authenticated public classical communication. Rates are per input
copy. Security includes a purification and the full public transcript.
States, bases and decompositions are known. $E_R$ denotes relative
entropy of entanglement and $E_R^\infty$ its regularization.
We use the established inequalities
\[
 E_D\le K_D\le E_R^\infty\le E_R,
\]
with the key upper bound from \cite[Theorem 9]{hhho}.
Logarithms are base two. Write
\[
 h(x)=-x\log_2x-(1-x)\log_2(1-x),
\]
with zero terms interpreted by continuity, and set
\begin{equation}\label{eq:b}
 b(x,y)=(x+y)h\!\left(\frac{y}{x+y}\right),\qquad b(0,0)=0.
\end{equation}

\begin{lemma}[A seed from a dominated component]\label{lem:seed}
Let $\rho=a\theta+(1-a)\chi$, where $0<a\le1$, and put
$u_A=S(\theta_A)-S(\theta)$. If $a2^{u_A}>1$, then $\rho$ has
a strictly positive two-way distillation rate.
\end{lemma}
\begin{proof}
Put $X=\rho_A$. If the reduction criterion held, then
$a\theta\le\rho\le X\otimes I$. Restrict to $\operatorname{supp}X$;
this contains the support of $\theta_A$, since $X\ge a\theta_A$.
Operator monotonicity of the logarithm, applied with a positive
regularization when $\theta$ is singular, gives
\[
 \log_2a-S(\theta)\le\Tr\theta_A\log_2X,
 \qquad
 \log_2a+u_A+D(\theta_A\Vert X)\le0.
\]
Nonnegativity of relative entropy contradicts the hypothesis.
Violation of reduction implies a positive finite-state distillation
rate \cite[Section VI]{hh}. This uses the established logarithmic
method relating reduction to conditional entropy
\cite[Theorem 2, equations (21)--(24)]{vwentropy}, rather than a new
general distillability criterion. The finite-dimensional filtering and
isotropic-twirling protocol supplies the seed. Its twirl has an exact
finite convex-unitary realization by Carath\'eodory's theorem, so no
uncountable outcome convention is needed. The seed rate may be small;
the state is fixed before the coding block length tends to infinity.
\end{proof}

\begin{theorem}[Stopped classification]\label{thm:classification}
For any state $\rho$ on $\mathbb C^d\otimes\mathbb C^d$ in a specified
paired product basis, let
\[
 B_{jk}=\langle jj|\rho|kk\rangle,\quad w=\Tr B,\quad
 v=w\{S(\operatorname{diag}(B/w))-S(B/w)\},
\]
setting $v=0$ when $w=0$. Then
\begin{equation}\label{eq:classlower}
 E_D(\rho)\ge\max\{0,v-h(w)-2(1-w)\}.
\end{equation}
If $\rho$ is invariant under every $U\otimes\overline U$ with $U$
diagonal unitary (the established CLDUI class), then
\begin{equation}\label{eq:cldui}
 \max\{0,v-h(w)-2(1-w)\}
 \le E_D(\rho)\le K_D(\rho)\le E_R^\infty(\rho)\le E_R(\rho)\le v.
\end{equation}
The interval has width at most $\log_2 5$, independently of $d$.
\end{theorem}
\begin{proof}
\emph{Finite phase averaging.} Apply shared phases
$U=\operatorname{diag}(\zeta^{a_j})$, $a_j\in\{0,1,2\}$ and
$\zeta=e^{2\pi i/3}$, as $U\otimes\overline U$.
With $\Pi_{\rm eq}=\sum_j|jj\rangle\langle jj|$, the averaged state is
\begin{equation}\label{eq:twirl}
 T(\rho)=\Pi_{\rm eq}\rho\Pi_{\rm eq}
       +\sum_{j\ne k}\langle jk|\rho|jk\rangle|jk\rangle\langle jk|.
\end{equation}
The characters $a_j-a_k$ are distinct for ordered pairs $j\ne k$,
whereas all $|jj\rangle$ have zero character. This finite LOCC map
preserves the entire equality block, including mixed maximally
correlated states. The averaged $AB$ marginal is block diagonal in
$z=j\mathbin\oplus k$ after padding labels to $m=\lceil\log_2d\rceil$
bits. This is not an assertion about a state conditioned on the public
phase record or on a purifying environment. The cases $d=1$ or $w=0$
have a trivial lower bound. When $w=1$, ordinary hashing gives $E_D\ge v$.

\emph{A fixed low-entropy classification.} For invertible binary
$m\times m$ matrices $A$, let $J_A$ be the first nonzero coordinate
of $Az$, or zero for $z=0$. Put $q=1-w$. Uniformly averaging $A$
sends each nonzero $z$ uniformly over nonzero vectors, so
\[
 \Pr(J=\ell\mid J>0)=\frac{2^{m-\ell}}{2^m-1},\qquad 1\le\ell\le m.
\]
Its entropy is $2-m/(2^m-1)+\log_2(1-2^{-m})\le2$.
Concavity gives $\mathbb E_A H(J_A)\le h(q)+2q$.
Thus one fixed $A$ satisfies this bound and can be used on all copies.
No efficient search for this matrix is assumed.

\emph{Coherent queries and their cost.} For binary local controls
$a(j),b(k)$, append a qubit Bell pair and apply controlled $X$ on
each party's half. Computational-basis outcomes $u,v$ on that pair
induce the source Kraus operator
\begin{equation}\label{eq:kraus}
 K_{uv}=2^{-1/2}\Pi^{\rm par}_{u\oplus v},\qquad
 \Pi^{\rm par}_y=\sum_{a(j)\oplus b(k)=y}|jk\rangle\langle jk|.
\end{equation}
Indeed $(X^{a(j)}\otimes X^{b(k)})|\Phi_2\rangle$
depends only on $a(j)\oplus b(k)$. The redundant measured bit reveals
no Schmidt index. The same formula holds for parities over subsets
of source copies. Each query consumes one ebit while preserving
coherences inside a parity eigenspace. This is the established
appended-pair/partial-breeding primitive \cite{hostens,vv}.

At stage $\ell$, learn the next bit only on copies with previous
bits zero, discarding a copy as soon as a one is learned. Let
$\pi_\ell=\Pr(J_A=\ell)$ and
$r_\ell=\Pr(J_A=0\text{ or }J_A\ge\ell)$. Binary typical-set hashing
of these copies costs $r_\ell h(\pi_\ell/r_\ell)$ ebits per original
copy, and the stopped-prefix chain rule gives
\begin{equation}\label{eq:prefix}
 C:=\sum_{\ell=1}^m r_\ell h(\pi_\ell/r_\ell)
   =H(J_A)\le h(w)+2(1-w).
\end{equation}
Zero-weight terms are omitted. For clarity, a typical set of binary
strings of length $k$ has at most $2^{k(H+\eta)}$ members; a random
linear hash with $k(H+2\eta)$ rows collides with another typical string
with probability at most $2^{-k\eta}$, and atypical probability vanishes.
If that row count exceeds $k$, measure all $k$ bits instead. Every row
is implemented by \eqref{eq:kraus}. Typical budgets cover random
sub-sample sizes, with abort on atypical counts. There are finitely
many stages at fixed $d$. On the averaged block-diagonal marginal,
correct classification acts exactly as the ideal block projectors;
the trace-norm error is at most twice the classification error.
No individual Schmidt label is measured.

\emph{Survivors and seed.} The all-zero projector is $\Pi_{\rm eq}$.
Survivors occur with probability $w$ and have state
$\theta=\Pi_{\rm eq}\rho\Pi_{\rm eq}/w$, whose correlated-basis matrix
is $B/w$. Write $u=S(\theta_A)-S(\theta)=v/w$; this is nonnegative
because dephasing increases entropy. Hashing \cite[Theorem 10]{dw}
achieves gross yield $G=wu=v$. Correct classification leaves independent
identical survivor states conditional on their positions. A fixed
output length below their typical count handles random numbers.

If $v-h(w)-2(1-w)\le0$, the asserted lower bound is trivial.
Otherwise $wu>h(w)\ge-w\log_2w$, so $w2^u>1$.
The averaged state is $w\theta+(1-w)\chi_{\rm diag}$ with a separable
mismatch state. Lemma~\ref{lem:seed} therefore supplies a positive
unassisted seed rate. This establishes the seed from the actual
survivor state, including when it is mixed.

\emph{Paying for and recycling the reserve.} Fix $\eta$ with
$0<2\eta<G-C$. A sufficiently long block of length $\ell$ consumes
at most $(C+\eta)\ell$ ebits and returns at least $(G-\eta)\ell$,
with joint error $\varepsilon_\ell\to0$. Generate the first reserve
from $O(\ell)$ source copies at the positive fixed seed rate. Run
$t_\ell\to\infty$ blocks, retaining a reserve from each output,
and choose $t_\ell$ slowly enough that $t_\ell\varepsilon_\ell\to0$.
The initial fraction is $O(1/t_\ell)$; contractivity and telescoping
bound the joint error by the seed error plus $t_\ell\varepsilon_\ell$.
Approximate reserves are included in that error. A diagonal choice of
lengths followed by $\eta\downarrow0$ achieves net rate $G-C$.
This standard removal of assistance \cite[Section V]{vv} needs no
hashing error exponent. Final fidelity to pure Bell pairs implies
decoupling from any purification and the complete public transcript.
No free entanglement or restricted Eve has been introduced.

For a CLDUI input, \eqref{eq:twirl} already equals $\rho$.
Its fully dephased state $\Delta\rho$ is separable, and the direct-sum
entropy formula gives
$D(\rho\Vert\Delta\rho)=S(\Delta\rho)-S(\rho)=v$.
Thus $E_R(\rho)\le v$, giving \eqref{eq:cldui}.
The width is bounded by $h(w)+2(1-w)\le\log_2 5$,
including when the lower endpoint is zero.
\end{proof}

\begin{corollary}[Diagonal orthogonal symmetry]\label{cor:ldoi}
Suppose $\rho$ is invariant under every $O\otimes O$ with $O$ a
diagonal sign matrix (the established LDOI class). Define $B,w,v$ as
above and $\mathcal C_\Delta(\rho)=S(\Delta\rho)-S(\rho)$.
Then
\begin{align}
 \max\{0,v-h(w)-2(1-w)\}&\le E_D(\rho)\le K_D(\rho)
 \le E_R^\infty(\rho)\le E_R(\rho)\le\mathcal C_\Delta(\rho),\label{eq:ldoi}\\
 0\le\mathcal C_\Delta(\rho)-E_D(\rho)&\le h(w)+3(1-w)\le\log_2 9.
 \label{eq:log9}
\end{align}
In particular the same constant bounds $E_R-E_D$ and $K_D-E_D$.
\end{corollary}
\begin{proof}
The known block characterization \cite[Proposition 2.2 and equations
(5)--(8)]{sn} consists of $B$ on the equality subspace and positive
$2\times2$ blocks $R_{ij}$ on $\operatorname{span}\{|ij\rangle,|ji\rangle\}$,
$i<j$. These subspaces are globally orthogonal; we do not measure
them by a presumed free nonlocal operation. Set $p_{ij}=\Tr R_{ij}$.
The entropy identity for this direct sum gives
\[
 \mathcal C_\Delta(\rho)=v+c,\qquad
 c=\sum_{i<j:p_{ij}>0}p_{ij}
 \{S(\operatorname{diag}(R_{ij}/p_{ij}))-S(R_{ij}/p_{ij})\}.
\]
Each bracket lies in $[0,1]$, hence $0\le c\le1-w$.
Theorem~\ref{thm:classification} gives the lower bound by the finite
local phase twirl, which preserves $B$ and removes the swap coherences.
The separable reference $\Delta\rho$ gives the upper bound.
The interval width is at most $h(w)+2(1-w)+c\le h(w)+3(1-w)$,
whose maximum is $\log_2 9$ at $w=1/9$.
\end{proof}
The bounds involve the spectrum of a $d\times d$ block and, for the
sharper upper endpoint, $2\times2$ blocks. Evaluating these quantities
is distinct from implementing an efficient distillation protocol.
Neither constant is claimed optimal. The upper endpoints apply to
the stated invariant classes; for an arbitrary state only
\eqref{eq:classlower} is asserted. No general bound on $E_C-E_D$
or solution of the NPT bound-entanglement problem follows.

\begin{theorem}[An arbitrary noisy maximally correlated component]\label{thm:recovery}
Let $\theta=\sum_{jk}t_{jk}|jj\rangle\langle kk|$ be a state,
$\sigma=\theta_A$, and $u=S(\sigma)-S(\theta)$. Let $\chi$ be any
state on the same $d\times d$ space, with no marginal condition.
For $\lambda,\delta\ge0$, $\lambda+\delta\le1$, define
\[
 \omega=\lambda\theta+\delta\chi+
   (1-\lambda-\delta)\sigma\otimes|e\rangle\langle e|,
\]
where $e$ is an orthogonal receiver flag. Then
\begin{equation}\label{eq:recovery}
 E_D(\omega)\ge\max\{0,\lambda u-b(\lambda,\delta)-\delta\log_2 5\}.
\end{equation}
If $\theta=\psi$ is pure, write $s=S(\psi_A)$. For separable $\chi$,
$E_C(\omega)\le\lambda s$ and
\begin{equation}\label{eq:irreversible}
 0\le E_C(\omega)-E_D(\omega)\le b(\lambda,\delta)+\delta\log_2 5.
\end{equation}
In particular, every mixture $\rho=\lambda\psi+(1-\lambda)\chi$ of
one pure state and an arbitrary separable state satisfies
\begin{equation}\label{eq:log6}
 E_C(\rho)-E_D(\rho)\le h(\lambda)+(1-\lambda)\log_2 5\le\log_2 6.
\end{equation}
\end{theorem}
\begin{proof}
Bob announces the code/flag outcome; put $P=\lambda+\delta$.
The case $P=0$ is trivial. Apply Theorem~\ref{thm:classification}
to the accepted normalized state. Put
$\kappa=\Tr\Pi_{\rm eq}\chi$, $P_0=\lambda+\delta\kappa$ and
$\Gamma=\lambda\theta+\delta\Pi_{\rm eq}\chi\Pi_{\rm eq}$.
When $\kappa>0$, the normalized projected noise $\xi$ is maximally
correlated and has $I_c(\xi)\ge0$. Entropy concavity and the
entropy-of-mixture upper bound yield
\[
 P_0 I_c(\Gamma/P_0)
 \ge\lambda u+\delta\kappa I_c(\xi)-b(\lambda,\delta\kappa)
 \ge\lambda u-b(\lambda,\delta\kappa).
\]
The physical rate is the accepted rate multiplied by $P$.
The classification cost is at most
$P h(\delta(1-\kappa)/P)+2\delta(1-\kappa)$.
The three-branch entropy identity gives
\[
 b(\lambda,\delta\kappa)+P h\!\left(\frac{\delta(1-\kappa)}P\right)
 =b(\lambda,\delta)+\delta h(\kappa).
\]
Since $h(\kappa)+2(1-\kappa)\le\log_2 5$, this proves
\eqref{eq:recovery}; zero-weight branches are interpreted directly.
The seed and reserve are already paid for by
Theorem~\ref{thm:classification}, which uses the actual survivor.
For pure $\theta$ and separable $\chi$, formation by typical branch
strings uses pure-state dilution \cite{bbps} at fraction $\lambda$
and free separable preparation otherwise, giving $E_C\le\lambda s$.
Combine with $E_D\le E_C$. The last maximum is $\log_2 6$ at
$\lambda=1/6$. These are upper bounds, not attained gaps established here.
\end{proof}

\begin{corollary}[Mixed good component and separable contamination]\label{cor:mixedkey}
For $\rho=\lambda\theta+(1-\lambda)\chi$ with maximally correlated
$\theta$ and separable $\chi$, each of $E_D,K_D,E_R^\infty,E_R$ lies
between $\max\{0,\lambda u-h(\lambda)-(1-\lambda)\log_2 5\}$ and
$\lambda u$. Consequently these rates converge absolutely to
$(1-q_n)u_n$ for $\lambda_n=1-q_n\to1$, even in growing dimensions.
\end{corollary}
\begin{proof}
The lower bound is \eqref{eq:recovery} with $\delta=1-\lambda$.
For $\theta$, its diagonal separable reference gives $E_R(\theta)\le u$.
Convexity of $E_R$ gives $E_R(\rho)\le\lambda u$, completing the
sandwich. No reversibility of $E_C$ is implied for a mixed $\theta$.
\end{proof}

\begin{corollary}\label{cor:separable}
Let $\rho_n=(1-q_n)\psi_n+q_n\chi_n$, where $\psi_n$ is pure and
$\chi_n$ is separable, in possibly growing finite dimensions. If
$q_n\to0$, then $E_C(\rho_n)-E_D(\rho_n)\to0$ without requiring
$q_n\log d_n\to0$. Each of $E_D,E_C$ lies between
$(1-q_n)S((\psi_n)_A)-h(q_n)-q_n\log_2 5$ and
$(1-q_n)S((\psi_n)_A)$.
\end{corollary}
For example, take $\psi_n=\Phi_{d_n}$,
$q_n=(\log_2d_n)^{-1/2}$, and $d_n\to\infty$.
The absolute irreversibility vanishes although $q_n\log_2d_n$ diverges,
for every choice of separable $\chi_n$.

\begin{corollary}\label{cor:ebchannel}
For any entanglement-breaking channel $\mathcal B_d$ from
$\mathbb C^d$ to itself, let
$\mathcal N_d=\lambda\id_d+(1-\lambda)\mathcal B_d$.
With free two-way classical communication and no initial entanglement
or key, each of $Q_2(\mathcal N_d)$ and $K_2(\mathcal N_d)$ lies in
\begin{equation}\label{eq:ebchannel}
 \left[\max\{0,\lambda\log_2d-h(\lambda)-(1-\lambda)\log_2 5\},
       \lambda\log_2d\right].
\end{equation}
The gap bound is uniform over $d$ and $\mathcal B_d$, and at most
$\log_2 6$. It tends to zero as $\lambda\to1$, with no relation
required between $1-\lambda$ and $d$.
\end{corollary}
\begin{proof}
Transmit halves of $\Phi_d$. The noisy Choi branch is separable with
Alice marginal $I/d$, so Theorem~\ref{thm:recovery} proves the lower
bound. Ebits yield quantum communication and secret key.
For any purified input, an extension that records the branch and a
separable decomposition of the entanglement-breaking output has
conditional mutual information $2\lambda S(\rho)$. Thus channel
squashed entanglement is at most $\lambda\log_2d$.
Its adaptive upper bound on $K_2$, and $Q_2\le K_2$, follow from
\cite{dsw} in finite dimension. The channel convexity and vanishing
property on entanglement-breaking channels are also established in
\cite[Sections III.B--C]{gew}. No covariance or reset form of
$\mathcal B_d$ is required.
\end{proof}

\section{Arbitrary noise and worst-case communication}\label{sec:minimax}
Let $U:\mathbb C^d\to\mathcal H_B$ be an isometry, $d\ge2$, with
finite-dimensional output, and consider the fixed memoryless channel
\[
 \mathcal N_{\mathcal M}=\lambda\operatorname{Ad}_U+(1-\lambda)\mathcal M,
 \qquad q=1-\lambda,
\]
where $\mathcal M$ is any CPTP map. It need not be separable, unital,
covariant or a replacement channel. We use unassisted classical,
quantum and private capacities $C,Q,P$, entanglement-assisted classical
capacity $C_{\rm EA}$, and two-way quantum and key capacities $Q_2,K_2$.
Only $C_{\rm EA}$ receives free initial entanglement; only $Q_2,K_2$
receive free authenticated public communication. Each block has finite
memories, alphabets and rounds, with no uniform bound across blocklengths.
All private tasks include the full channel environment and public record.
The infima below optimize over fixed channels, with codes allowed to
depend on the channel. They are not compound-channel or adversarial-noise
capacities.

\begin{corollary}[Six noise-universal bounds]\label{cor:universal}
Define the benchmarks and losses by
\[
\begin{array}{c|c|c}
 X&r_X&c_X\\ \hline
 C&\lambda\log_2d&h(\lambda)\\
 C_{\rm EA}&2\lambda\log_2d&h(\lambda)\\
 Q,P&(2\lambda-1)_+\log_2d&h(\lambda)\\
 Q_2,K_2&\lambda\log_2d&h(\lambda)+q\log_2 5
\end{array}
\]
For every $\mathcal M$, $X(\mathcal N_{\mathcal M})\ge
\max\{0,r_X-c_X\}$, and
\begin{equation}\label{eq:minimaxbounds}
 \max\{0,r_X-c_X\}\le\inf_{\mathcal M}X(\mathcal N_{\mathcal M})\le r_X.
\end{equation}
For $\lambda\le1/2$, the $Q$ and $P$ infima are exactly zero.
The losses are at most one bit for the first four capacities and
$\log_2 6$ for the last two; all tend to zero as $\lambda\to1$,
uniformly in dimension and noise.
\end{corollary}
\begin{proof}
The first four lower bounds are standard entropy-of-mixture estimates.
For a uniform orthonormal input ensemble, put
$\omega_i=\mathcal M(|i\rangle\langle i|)$ and
$\bar\omega=d^{-1}\sum_i\omega_i$. Concavity and the mixture upper bound
give Holevo information at least
$\lambda\log_2d-h(\lambda)+q[S(\bar\omega)-d^{-1}\sum_iS(\omega_i)]$.
For the Bell input write the output as
$\rho=\lambda\Phi_U+q\chi$, where $\chi_A=I/d$. Then
\[
 S(\rho_B)\ge\lambda\log_2d+qS(\chi_B),\qquad
 S(\rho_{AB})\le h(\lambda)+q[\log_2d+S(\chi_B)].
\]
Coherent-information coding and $P\ge Q$ give the $Q,P$ bounds;
adding $S(\rho_A)=\log_2d$ and the entanglement-assisted coding theorem
gives the $C_{\rm EA}$ bound \cite{wq,bsst}.

For $Q_2,K_2$, Bob projects onto $\operatorname{ran}U$. The retained
unnormalized state is $\lambda\Phi_d+\delta\chi'$ for some
$0\le\delta\le q$. Apply Theorem~\ref{thm:recovery} to its normalized
state and multiply the achievable rate by $\lambda+\delta$.
This gives $\lambda\log_2d-b(\lambda,\delta)-\delta\log_2 5$,
at least the asserted bound. The rejected state's Alice marginal is
irrelevant: those copies are discarded. Distributing Choi copies and
distilling them is allowed for $Q_2$, and $K_2\ge Q_2$.

For the upper bound on the infimum, choose any replacement noise.
The resulting channel is receiver postprocessing of the erasure channel
with survival $\lambda$, whose six capacities are $r_X$
\cite{bds,bsst,dsw}. For completeness, erasure has coherent-information
rate $(2\lambda-1)\log_2d$ when degradable and has $Q=P=0$ when
antidegradable; its two-way quantum/key rate is $\lambda\log_2d$,
achieved on successful transmissions and bounded by the erasure
squashed-entanglement bound. This upper bound concerns the infimum,
not every noise: $\mathcal M=\operatorname{Ad}_U$ is noiseless.
\end{proof}

The worst noise admits a standard symmetry reduction. We give the
resource accounting because private and unassisted capacities must
not silently receive free shared secret randomness. Define the
depolarizing channel $\mathcal D_f$ through its normalized Choi state
\begin{equation}\label{eq:isotropic}
 J_f=f\Phi_d+\frac{1-f}{d^2-1}(I-\Phi_d),\qquad
 \mathcal D_f(\rho)=p_f\rho+(1-p_f)\frac{I}{d}\Tr\rho,
 \quad p_f=\frac{d^2f-1}{d^2-1}.
\end{equation}

\begin{proposition}[Simultaneous minimax reduction]\label{prop:minimax}
For each of the six capacities above,
\begin{equation}\label{eq:exactminimax}
 \inf_{\mathcal M}X(\mathcal N_{\mathcal M})
 =X(\mathcal D_{f_*}),\qquad f_* = \max\{\lambda,d^{-2}\}.
\end{equation}
One noise channel attains all six infima simultaneously. This reduces
the problem to depolarizing capacities; it does not evaluate the
generally unknown quantum, private or two-way capacities.
\end{proposition}
\begin{proof}
First suppose $U=I$. Conjugation twirling sends any channel $\mathcal N$
to $\mathcal D_f$, where $f=\langle\Phi_d|J_{\mathcal N}|\Phi_d\rangle$.
The Haar average is a finite convex combination of unitary conjugations
by finite-dimensional convexity. Twirling cannot increase any of the
six capacities. For $C,Q,C_{\rm EA}$, expand an $n$-use code for the
twirled channel over the independent unitary seeds. Average decoding
error or entanglement infidelity tends to zero, so one deterministic
sequence works; absorb it into the encoder and decoder. This is the
usual coding use of twirling \cite[Appendix B]{bdsw}.

For $P$, a Stinespring dilation of the twirl retains a coherent seed
register. Dephasing that register in Eve's output gives the actual
environment of $\mathcal N$ with the public seed. Contractivity gives
security for this degraded environment. Since the seed is independent
of the message, the block diagonal secrecy distance is the average
conditional secrecy distance. A deterministic sequence has vanishing
sum of average decoding and secrecy errors. Thus no random resource
is added to the unassisted code. For adaptive $Q_2,K_2$, fresh public
seeds implement each twirled use. Measuring the coherent environment
seed gives precisely this implementation, with every seed retained by
Eve. The original protocol ignores those extra labels after the required
conjugations; the full output security is preserved by contractivity.

For $\mathcal N=\lambda\id+q\mathcal M$ we have $f\ge\lambda$.
When $\lambda\ge d^{-2}$, $p_f\ge p_\lambda\ge0$; receiver
postprocessing by a depolarizing channel of parameter $p_\lambda/p_f$
reduces $\mathcal D_f$ to $\mathcal D_\lambda$. The case $p_f=0$ is
immediate. Hence $X(\mathcal N)\ge X(\mathcal D_\lambda)$.
Equality is attained by the CPTP noise
\[
 \mathcal M_0(\rho)=\frac{dI\Tr\rho-\rho}{d^2-1},
 \qquad J_{\mathcal M_0}=\frac{I-\Phi_d}{d^2-1}.
\]
When $\lambda\le d^{-2}$, the CPTP map
$\mathcal M=(\mathcal D_{d^{-2}}-\lambda\id)/(1-\lambda)$
makes the mixture completely depolarizing. Its Choi eigenvalues are
$(d^{-2}-\lambda)/(1-\lambda)$ on $\Phi_d$ and
$[d^2(1-\lambda)]^{-1}$ on its orthogonal complement.
All six capacities then vanish.

For general $U$, decode by
$\mathcal R(Y)=U^*YU+\Tr[(I-UU^*)Y]\tau_0$, with any state $\tau_0$.
This CPTP map satisfies $\mathcal R\operatorname{Ad}_U=\id$,
so the lower bound reduces to the square case. Composing its attaining
noises with $\operatorname{Ad}_U$ gives equality.
\end{proof}

The classical formulas in this reduction are established results.
For $\lambda\ge d^{-2}$, put $a=(d\lambda+1)/(d+1)$. King
\cite[Theorem 1]{king} and Bennett et al.\ \cite{bsst} give, respectively,
\begin{align*}
 \inf_{\mathcal M}C(\mathcal N_{\mathcal M})
 &=\log_2d-h(a)-(1-a)\log_2(d-1),\\
 \inf_{\mathcal M}C_{\rm EA}(\mathcal N_{\mathcal M})
 &=2\log_2d-h(\lambda)-(1-\lambda)\log_2(d^2-1).
\end{align*}
Below $d^{-2}$ both are zero. The new classification bound instead
controls the unresolved two-way rates with an absolute interval.

\begin{corollary}[An absolute isotropic interval]\label{cor:isotropic}
For $f\ge1/d$, set
\[
 q=1-f,\quad t=\frac{q}{d+1},\quad w=f+t,\quad
 v=w\log_2d-wh(f/w)-t\log_2(d-1),
\]
and define
\[
 L=\max\{0,v-h(w)-2(1-w)\},\qquad
 U=\log_2d-h(f)-q\log_2(d-1).
\]
Both $E_D(J_f)$ and $K_D(J_f)$ lie in $[L,U]$, whose width is exactly
\begin{equation}\label{eq:isowidth}
 U-L=\min\{U,qA_d\},\qquad
 A_d=\frac{2d}{d+1}+\log_2(d+1)-\frac{d}{d+1}\log_2(d-1).
\end{equation}
The same interval bounds the two minimax capacities in
\eqref{eq:exactminimax} with $f=\lambda\ge1/d$.
For fixed $0<f<1$, its width tends to $2(1-f)$ as $d\to\infty$.
The two-way minimax capacities vanish for $\lambda\le1/d$ and are
strictly positive for $\lambda>1/d$.
\end{corollary}
\begin{proof}
The correlated block of $J_f$ has eigenvalues $f$ and
$q/(d^2-1)$ with multiplicity $d-1$, total weight $w$ and uniform
diagonal. Thus Theorem~\ref{thm:classification} gives $L$.
The known isotropic relative entropy is $E_R(J_f)=U$
\cite[Supplementary Note 5, Eqs.\ (S195)--(S196)]{plob}; it bounds
both rates. The entropy chain rule
\[
 h(w)+wh(f/w)=h(f)+q h\!\left(\frac1{d+1}\right)
\]
gives $U-[v-h(w)-2(1-w)]=qA_d$, proving \eqref{eq:isowidth}.
Depolarizing channels are teleportation covariant, so their two-way
capacities equal $E_D(J_f)$ and $K_D(J_f)$ respectively: distribution
gives one direction and exact teleportation simulation gives the other
\cite{bdsw,plob}. Isotropic states are separable up to $f=1/d$;
above it, reduction violation gives positive distillability
\cite{hh,vw}. This proves the threshold, even where the displayed
quantitative lower bound is zero.
\end{proof}
Vollbrecht--Wolf \cite[Section VI.B--C, Eqs.\ (40)--(47)]{vw}
already obtained the optimal normalized isotropic rate
$E_D(J_f)/\log_2d\to f$. The assertion here is the finite-dimensional
absolute interval \eqref{eq:isowidth}, not that prelog, the symmetry
reduction, or the established classical and entanglement-assisted
formulas. No equality $E_D=K_D$, optimal constant, or exact positive
depolarizing quantum/private capacity is claimed.


\section{Oscillator model, prior work and conventions}
Let $\mathcal H=\ell^2(\mathbb N_0)$ and
$H=\sum_{j\ge0}j|j\rangle\langle j|$. Fix $0\le\lambda\le1$ and a density
operator $\tau$, initially without energy or entropy restrictions.
When finite, write $M=\Tr H\tau$. The memoryless reset channel is
\begin{equation}\label{eq:reset}
 \Rcal(\rho)=\lambda\rho+(1-\lambda)\tau\Tr\rho.
\end{equation}
Here $\lambda$ is a survival probability, not a beam-splitter transmissivity.
This is not the Gaussian thermal-loss channel. No assumption $[H,\tau]=0$
is made. The comparison erasure channel is
\[
 \Ecal(\rho)=\lambda\rho\oplus(1-\lambda)\Tr\rho\,|e\rangle\langle e|,
\]
with an additional orthogonal flag dimension.
Entropies are in bits, except for an explicit $\ln$. Put
\[
 g(N)=(N+1)\log_2(N+1)-N\log_2N,
\]
interpreting zero terms by continuity. The Gibbs variational principle gives
$\sup_{\Tr H\rho\le N}S(\rho)=g(N)$, attained by the geometric distribution.

Bennett, DiVincenzo and Smolin established finite-dimensional erasure
capacities \cite{bds}. Zhong, Oh and Jiang studied continuous-variable
erasure with an external flag, explicitly distinguishing it from
replacement by the vacuum without a flag \cite{zoj}.
Catalano et al.\ recently studied mixtures of bosonic pure-loss channels
\cite{catalano}. Their Section IV, equations (37)--(42), treats exactly
the identity/vacuum-reset case and optimizes its entanglement-assisted
capacity. Section VI.C, equation (48), gives coherent information
$(2\lambda-1)S(\rho)$ on inputs orthogonal to the vacuum.
That formula, together with the erasure upper bound, already yields
the vacuum case of our quantum asymptotic statement.
Their structural results also include non-entanglement-breaking behavior
of nontrivial fading channels. These special cases are not claimed as new.
Strong classical additivity of the finite-dimensional erasure channel
is also established, for example by Amosov--Mancini
\cite[Theorem 9]{amosov}; our partner statement transfers an erasure
upper bound to reset with a vanishing defect.
Catalano et al.\ also study unconstrained two-way capacities of finite
pure-loss mixtures \cite[Section III.D, Theorem 1]{catalano}, including
their divergence in the presence of an identity component.
We do not claim that divergence or the vacuum-reset mechanism as new.

The hidden-branch estimate below proves four absolute-gap limits
without reset energy or entropy assumptions. Theorem~\ref{thm:recovery}
then gives two further assisted limits and state-entanglement
reversibility under the same unrestricted fixed-state assumption.
The coding, hashing and pure-state dilution dependencies are established
\cite{holevo,hs,wq,dw,dsw,bbps}.

\paragraph{Operational conventions.}
For $n$ uses let $H^{(n)}=\sum_i H_i$. The input allowance is mean
total energy at most $nN$, averaged over encoded inputs.
$C(\mathcal N,N)$ is classical capacity with vanishing decoding error,
no prior entanglement and no feedback.
$C_{\rm EA}(\mathcal N,N)$ allows unlimited prior entanglement,
with the cost on channel inputs.
$Q(\mathcal N,N)$ is entanglement-transmission capacity with the cost
on the encoded maximally mixed logical state.
$P(\mathcal N,N)$ is unassisted private classical transmission, with
no preshared key or free public discussion.
Messages are equiprobable; decoding error is averaged over messages.
For $Q$ the entanglement fidelity tends to one. For $P$ the average
trace distance of each full-environment state from a common fixed state
tends to zero. The security environment is that of a full dilation of
the channel, not just the mathematical reset-branch register.
Wilde--Qi's operational equivalences \cite[Theorem 1]{wq}
apply to a general channel and energy observable; they identify these
quantum and private conventions with the corresponding uniform-energy
quantum and secret-key transmission tasks. The finite-output-entropy
coding formulas in their Theorems 2--3 will be used only on the finite
decoded channel or the erasure comparison channel.
Their secret-key transmission task does not allow free two-way communication.
The allowance is across the block; there is no peak constraint on each
individual use. We consider vanishing-error capacities, not
strong-converse rates at a fixed nonzero error.

\begin{theorem}\label{thm:main}
For every fixed $\lambda$ and $\tau$ as above, as $N\to\infty$,
\begin{align}
 C(\Rcal,N)&=\lambda g(N)+o(1),\label{eq:c}\\
 C_{\rm EA}(\Rcal,N)&=2\lambda g(N)+o(1),\label{eq:ea}\\
 Q(\Rcal,N)&=(2\lambda-1)_+g(N)+o(1),\label{eq:q}\\
 P(\Rcal,N)&=(2\lambda-1)_+g(N)+o(1).\label{eq:p}
\end{align}
Each leading term is an upper bound. The remainders for $Q$ and $P$
can differ; no equality at finite energy is claimed.
For $\lambda\le1/2$, $Q=P=0$ exactly for every $N$.
No finiteness of $S(\tau)$ or $\Tr H\tau$ is required.
The convergence is uniform over $\tau$ with $\Tr H\tau\le M_0$
for each fixed finite $M_0$, but no uniformity over all states is claimed.
\end{theorem}
The absolute error is stronger than matching only a coefficient of
$\log N$. One family of finite-rank inputs proves all four limits.
We obtain achievable rates through coding theorems, not an efficient
finite-block quantum coding algorithm.

\section{An exact erasure approximation}
For integers $K\ge1,d\ge2$, let $V|j\rangle=|K+j\rangle$
for $0\le j<d$, and $\Pi=VV^*$. Define the deterministic decoder
\begin{equation}\label{eq:decoder}
 D(X)=V^*\Pi X\Pi V\oplus\Tr[(I-\Pi)X]|e\rangle\langle e|.
\end{equation}
It measures membership in the code subspace, preserving quantum information
inside it, without postselecting away failed uses. Set
\begin{equation}\label{eq:delta}
 \tau_\Pi=V^*\tau V,\quad \epsilon=\Tr\Pi\tau,\quad
 \delta=(1-\lambda)\epsilon,\quad \mathcal V(\rho)=V\rho V^*.
\end{equation}
Denote the $d$-dimensional erasure channel by $\mathcal E_\lambda^{(d)}$.

\begin{lemma}\label{lem:diamond}
The decoded channel satisfies
\[
 D\Rcal\mathcal V(\rho)=
 [\lambda\rho+(1-\lambda)\tau_\Pi\Tr\rho]
 \oplus(1-\lambda)(1-\epsilon)\Tr\rho\,|e\rangle\langle e|,
\]
and
\begin{equation}\label{eq:diamond}
 \tfrac12\|D\Rcal\mathcal V-\mathcal E_\lambda^{(d)}\|_\diamond
 =\delta,\qquad
 \epsilon\le\Tr[\mathbf1_{\{H\ge K\}}\tau].
\end{equation}
\end{lemma}
\begin{proof}
Substitute \eqref{eq:reset} in \eqref{eq:decoder}. The difference from
erasure is $X\mapsto\Tr(X)A$, where
$A=(1-\lambda)(\tau_\Pi\oplus-\epsilon|e\rangle\langle e|)$.
This replacement map has diamond norm $\|A\|_1$:
partial trace has completely bounded trace norm one and a density
operator attains the norm. The positive and negative blocks have equal
absolute trace, so $\|A\|_1=2\delta$.
Since $\Pi\le\mathbf1_{\{H\ge K\}}$, the tail bound follows.
When $M<\infty$, $H\ge K\Pi$ also gives $\epsilon\le M/K$.
\end{proof}
The distance includes arbitrary references.
Reset coherences between $\Pi$ and its complement are removed by $D$,
so no square-root loss from gentle measurement is needed.

\section{A dimension-free information-loss bound}
Recall $b(\lambda,\delta)$ from \eqref{eq:b}, for
$0\le\delta\le1-\lambda$. It is the remaining uncertainty about a
binary branch after observing whether the output is in the code subspace.

\begin{lemma}\label{lem:hidden}
Let $\sigma$ be a $d$-dimensional state of entropy $s$, let $\psi_{RB}$
purify it, and let $\omega_{RB}=(\id\otimes D\Rcal\mathcal V)(\psi)$.
Then
\begin{equation}\label{eq:information}
 I(R:B)_\omega\ge2\lambda s-b(\lambda,\delta),\qquad
 I_c(\sigma,D\Rcal\mathcal V)\ge(2\lambda-1)s-b(\lambda,\delta).
\end{equation}
For a diagonal input ensemble and a number measurement after $D$,
the classical mutual information is at least
$\lambda s-b(\lambda,\delta)$.
\end{lemma}
\begin{proof}
Introduce a classical register $F$ recording transmit or reset:
\[
 \omega_{RBF}
 =\lambda\psi_{RB}\otimes|T\rangle\langle T|
 +(1-\lambda)\sigma_R\otimes D(\tau)_B\otimes|Z\rangle\langle Z|.
\]
The reference is independent of $F$, so $I(R:BF)=2\lambda s$.
The chain rule gives
$I(R:B)=2\lambda s-I(R:F|B)$.
Since $F$ is classical, $S(F|RB)\ge0$, and hence
$I(R:F|B)\le S(F|B)$.
Measure only code membership at $B$, producing the binary variable $Y$.
Data processing gives $S(F|B)\le H(F|Y)$.
An erasure output identifies reset; an accepted output has probability
$p=\lambda+\delta$ and conditional reset probability $q=\delta/p$.
Thus $H(F|Y)=p h(q)=b(\lambda,\delta)$.
Finally $I_c=I(R:B)-S(R)$.
For classical symbols $X$, the branch register is again independent
of the input and $I(X:Y_{\rm full},F)=\lambda s$.
The same chain rule, using membership as a coarsening of the full
number outcome, proves the classical bound.
All registers in this proof are finite dimensional.
\end{proof}
The bound does not scale with $\log d$. In particular,
$0\le b(\lambda,\delta)\le h(\delta)$: the rule that guesses
transmission on acceptance has error probability $\delta$, so its error
bit bounds $H(F|Y)$. The function is nondecreasing in $\delta$, since
$\partial_\delta b=\log_2(1+\lambda/\delta)$ when $\delta>0$.
Ordinary conditional-entropy continuity \cite{winter} would instead
introduce a term proportional to $\delta\log d$.

\begin{proposition}\label{prop:lower}
For a diagonal $\sigma=\sum_{j=0}^{d-1}p_j|j\rangle\langle j|$ of
entropy $s$, suppose $K+\sum_jjp_j<N$. Then
\begin{align}
 C(\Rcal,N)&\ge\lambda s-b(\lambda,\delta),\label{eq:lowerc}\\
 Q(\Rcal,N)&\ge\max\{0,(2\lambda-1)s-b(\lambda,\delta)\},\label{eq:lowerq}\\
 P(\Rcal,N)&\ge Q(\Rcal,N),\label{eq:lowerp}\\
 C_{\rm EA}(\Rcal,N)&\ge2\lambda s-b(\lambda,\delta).\label{eq:lowerea}
\end{align}
\end{proposition}
\begin{proof}
Apply Lemma~\ref{lem:hidden} and the coding theorems with strict cost
slack to the finite input and finite decoded output, with bounded cost
$V^*HV=\operatorname{diag}(K,\ldots,K+d-1)$.
The finite-alphabet classical coding theorem proves \eqref{eq:lowerc};
coherent-information achievability \cite[Theorem 2]{wq} proves
\eqref{eq:lowerq}; and the entanglement-assisted theorem
\cite[Theorem 1]{hs} proves \eqref{eq:lowerea}.
A quantum code yields a private code.
Every such code is implemented on the physical channel by $V$ and $D$.
This proves the lower bounds even when the physical reset output has
infinite entropy, without invoking a finite-output-entropy formula on
that channel itself.
\end{proof}

\begin{proposition}\label{prop:upper}
For every $N\ge0$,
\begin{align}
 C(\Rcal,N)&\le\lambda g(N),&
 C_{\rm EA}(\Rcal,N)&\le2\lambda g(N),\label{eq:upperc}\\
 Q(\Rcal,N)&\le(2\lambda-1)_+g(N),&
 P(\Rcal,N)&\le(2\lambda-1)_+g(N).\label{eq:upperq}
\end{align}
\end{proposition}
\begin{proof}
Reset is a postprocessing of $\Ecal$: keep the transmitted block and
replace the orthogonal flag by $\tau$. Every capacity is therefore
bounded by that of $\Ecal$ with the same input cost.

For a classical $n$-use ensemble, condition on the independent
surviving subset $J$. The Holevo information in that branch is at most
$\sum_{i\in J}S(\bar\rho_i)$. Its average is at most
$\lambda\sum_iS(\bar\rho_i)\le n\lambda g(N)$.
Geometric number ensembles approach equality.
For entanglement assistance,
$I(\rho,\Ecal)=2\lambda S(\rho)$; apply \cite[Theorem 1]{hs}.

For $\lambda\ge1/2$, erasure is degradable and has coherent information
$(2\lambda-1)S(\rho)$. Wilde--Qi's Theorem 3 \cite{wq} makes its
energy-constrained quantum and private capacities equal to optimized
coherent information. The Gibbs condition holds since
$\Tr e^{-\beta H}=(1-e^{-\beta})^{-1}<\infty$ for $\beta>0$.
Erasure output entropy is bounded by $h(\lambda)+\lambda g(N)$.
This verifies the finite-output-entropy condition on erasure, the
channel to which that theorem is applied. Its complementary output
entropies are finite by purification and subadditivity.
No such condition is required of the postprocessed reset channel. For $\lambda\le1/2$, erasure is antidegradable:
the environment can simulate the receiver, forcing its unassisted
private and quantum capacities to zero. Postprocessing preserves this
conclusion. Both arguments agree at $\lambda=1/2$.
\end{proof}

\section{Entropy-preserving escape}
A uniform distribution on a cutoff subspace loses a nonvanishing
constant relative to $g(N)$. Instead, for sufficiently large $N$ choose
\begin{equation}\label{eq:choose}
 \begin{gathered}
 K=\lceil\sqrt N\rceil,\quad m=N-K>0,\quad r=\frac{m}{m+1},\\
 d=\left\lceil2(m+1)\ln(N+1)\right\rceil,\quad t=r^d,\qquad
 p_j=\frac{(1-r)r^j}{1-r^d}\quad(0\le j<d).
 \end{gathered}
\end{equation}

\begin{lemma}\label{lem:geom}
For this distribution,
\begin{align}
 K+\sum_jjp_j&=K+m-\frac{dt}{1-t}<N,\label{eq:mean}\\
 s_N&=g(m)+\log_2(1-t)+\frac{dt}{1-t}\log_2r.\label{eq:entropy}
\end{align}
Also $t\le(N+1)^{-2}$ and
$a_N:=g(N)-s_N=O(N^{-1/2})+O(\log N/N^2)=o(1)$.
\end{lemma}
\begin{proof}
Differentiate the finite geometric sum to obtain the mean.
Substitute it into
$-\sum_jp_j[\log_2(1-r)-\log_2(1-t)+j\log_2r]$
to obtain the entropy. Since
$\ln r=-\ln(1+1/m)\le-1/(m+1)$, the stated tail bound follows.
The shift loss is at most
\[
 g(N)-g(N-K)=\int_{N-K}^N\log_2(1+1/x)\,dx
 \le K\log_2(1+1/(N-K))=O(N^{-1/2}).
\]
Finally $d|\ln r|=O(\ln N)$, so the other entropy terms have
absolute sum $O(t(1+d|\ln r|))$.
\end{proof}

\begin{proof}[Proof of Theorem~\ref{thm:main}]
With $\delta_N=(1-\lambda)\Tr\Pi_N\tau$, the preceding propositions give
\begin{align}
0\le\lambda g(N)-C(\Rcal,N)&\le\lambda a_N+b(\lambda,\delta_N),\label{eq:gapc}\\
0\le2\lambda g(N)-C_{\rm EA}(\Rcal,N)&\le2\lambda a_N+b(\lambda,\delta_N).\label{eq:gapea}
\end{align}
For $\lambda>1/2$, the gap of either $Q$ or $P$ below the common
erasure upper bound is at most
\begin{equation}\label{eq:gapq}
 (2\lambda-1)a_N+b(\lambda,\delta_N).
\end{equation}
Since $\tau$ is trace class and $K\to\infty$,
$\delta_N\le(1-\lambda)\Tr[\mathbf1_{\{H\ge K\}}\tau]\to0$.
The binary loss $b(\lambda,\delta_N)$ therefore vanishes without a
dimension-dependent multiplier.
If $M<\infty$, then $\delta_N\le(1-\lambda)M/K=O(M/\sqrt N)$.
Monotonicity of $b$ permits the replacement
$\delta_N\mapsto(1-\lambda)\min\{1,M/K\}$.
The resulting bounds are $O_M(\log N/\sqrt N)$ for fixed positive $M$;
when $M=0$ the branch loss is zero. The constants are uniform for $M\le M_0$.
For $\lambda\le1/2$ the quantum and private claims are already exact.
The endpoints $\lambda=0,1$ are the replacer and identity channels.
\end{proof}

Specific reset tails improve the rate. For thermal $\tau$ of mean $M>0$,
$\Tr\Pi\tau\le[M/(M+1)]^K$.
Taking $K=\lceil c\ln(N+1)\rceil$ with $c>2/\ln(1+1/M)$,
and the other choices in \eqref{eq:choose} with $m=N-K$, gives
all four gaps $O_M(\log N/N)$; \eqref{eq:tworate} below
also gives this rate for both two-way gaps.
This follows from the same estimates; optimal rates are not asserted.

For each fixed $N$ we first take coding block length to infinity.
Only afterwards do we take $N\to\infty$; no coding theorem uniform in
the growing dimension $d(N)$ is required.

\section{Two-way reset capacities and reversible distributed states}
$Q_2$ and $K_2$ for the oscillator use the countably decomposable LOCC
convention of \cite{dsw}, with free authenticated public two-way
classical discussion and no initial entanglement or secret key.
For adaptive $n$-use protocols the energy constraint is
$n^{-1}\sum_i\Tr H\rho^{(i)}_{\rm input}\le N$, averaging each
marginal over previous outcomes \cite[equation (101)]{dsw}.
Security includes the full channel environment and all public messages.

\begin{theorem}\label{thm:two}
For every fixed oscillator density operator $\tau$, with no entropy
or energy restriction,
\begin{equation}\label{eq:two}
 Q_2(\Rcal,N)=\lambda g(N)+o(1),\qquad
 K_2(\Rcal,N)=\lambda g(N)+o(1).
\end{equation}
Each leading term is an upper bound; the remainders need not agree.
For every finite encoding in Proposition~\ref{prop:lower},
\begin{equation}\label{eq:twofinite}
 Q_2(\Rcal,N)\ge
 \max\{0,\lambda s-b(\lambda,\delta)-\delta\log_2 5\},\qquad K_2\ge Q_2.
\end{equation}
\end{theorem}
\begin{proof}
For the upper bound use reset as receiver postprocessing of erasure.
Davis--Shirokov--Wilde \cite[Theorem 4 and equation (141)]{dsw} bound
both two-way capacities of erasure by its energy-constrained squashed
entanglement. Its Gibbs and finite-output-entropy hypotheses were
verified in Proposition~\ref{prop:upper}. For a purified input $\rho$,
retaining the erasure branch in a classical extension gives conditional
mutual information $2\lambda S(\rho)$; hence this squashed entanglement
is at most $\lambda g(N)$. This covers adaptive average-cost protocols
and is applied to erasure even when reset has infinite output entropy.

For the lower bound, transmit halves of the finite input purification
and apply $D$. The resulting state has the form of
Theorem~\ref{thm:recovery}, with
$\chi=\sigma_A\otimes(\tau_\Pi/\epsilon)$ if $\epsilon>0$,
and a pure useful component with $u=s$. This yields \eqref{eq:twofinite}; when
$\epsilon=0$ it is exactly flagged erasure. The seed also comes from
copies of this same finite-energy input, so no uncounted input cost
or initial entanglement is introduced.
For \eqref{eq:choose}, $s_N=g(N)-a_N$, $a_N\to0$ and
$\delta_N\to0$ for every fixed trace-class $\tau$.
Both loss terms in \eqref{eq:twofinite} vanish independently of its
entropy, proving the two limits separately. The endpoints $\lambda=0,1$
are respectively replacement and identity.
\end{proof}
For $M<\infty$, put $\bar\epsilon_N=\min\{1,M/K\}$ and
$\bar\delta_N=(1-\lambda)\bar\epsilon_N$. Each gap is bounded by
\begin{equation}\label{eq:tworate}
 \lambda a_N+b(\lambda,\bar\delta_N)+\bar\delta_N\log_2 5,
\end{equation}
since $b$ is increasing in its second argument. This is
$O_{M_0}(\log N/\sqrt N)$ uniformly for $M\le M_0<\infty$.
No rate uniform over all reset states is claimed.

\begin{corollary}\label{cor:state}
Let $\psi_N$ purify the encoded input in \eqref{eq:choose} and
$\Omega_N=(\id\otimes\Rcal)(\psi_N)$. For every fixed $\tau$,
\[
 E_D(\Omega_N)=\lambda g(N)+o(1),\qquad
 E_C(\Omega_N)=\lambda g(N)+o(1),\qquad
 E_C(\Omega_N)-E_D(\Omega_N)=o(1).
\]
The same holds after applying $D$ locally to the receiver.
\end{corollary}
\begin{proof}
Local decoding and Theorem~\ref{thm:recovery} give
$E_D(\Omega_N)\ge\lambda s_N-b(\lambda,\delta_N)-\delta_N\log_2 5$.
The decomposition
$\Omega_N=\lambda\psi_N+(1-\lambda)\sigma_A\otimes\tau$
gives $E_C(\Omega_N)\le\lambda s_N$ by the same random-branch
preparation and pure-state dilution. A density on a separable Hilbert
space has a countable spectral decomposition, so local preparation of
$\tau$ is allowed even at infinite entropy or energy. Decoding preserves
this upper bound. Use $E_D\le E_C$ and let $N\to\infty$.
\end{proof}
Local operations and preparation are not energy constrained in this
LOCC cost. This is a statement about the specified distributed states,
not the cost of simulating the channel. For each $N$ the operational
block length tends to infinity first. No finite-$N$ reversibility is asserted.

\section{Uniform classical additivity defect}
For finite-dimensional $\Psi$, let $C(\Rcal\otimes\Psi,N)$ mean $n$
parallel pairs, with mean oscillator energy at most $N$ per pair and
no extra constraint on the finite-dimensional input.

\begin{theorem}\label{thm:partner}
For every finite-dimensional $\Psi$,
\begin{align*}
0&\le C(\Rcal\otimes\Psi,N)-C(\Rcal,N)-C(\Psi)\\
 &\le\lambda g(N)-C(\Rcal,N).
\end{align*}
The supremum of this defect over all finite-dimensional partners tends
to zero as $N\to\infty$.
\end{theorem}
\begin{proof}
It suffices to show
$C(\Ecal\otimes\Psi,N)\le\lambda g(N)+C(\Psi)$.
For finite-dimensional $\mathcal M$, define its convex output-entropy roof
\[
 \widehat S_{\mathcal M}(\omega)=
 \inf_{\omega=\sum_xp_x|\psi_x\rangle\langle\psi_x|}
       \sum_xp_xS(\mathcal M(|\psi_x\rangle\langle\psi_x|)).
\]
If $B$ is finite dimensional and $A$ is possibly infinite dimensional,
\begin{equation}\label{eq:roof}
 S((\id_A\otimes\mathcal M)(\omega_{AB}))
       \ge\widehat S_{\mathcal M}(\omega_B).
\end{equation}
For pure $\omega_{AB}$, the left side is
$S(\mathcal M^c(\omega_B))$. Concavity on any pure-state decomposition
of $\omega_B$, and equality of output and complementary entropies on
pure inputs, give the lower bound.
For mixed $\omega_{AB}$, use its countable spectral decomposition,
entropy concavity and convexity of the finite-dimensional roof.
The application below has finite entropies by the energy bound.

Fix $n$ and $\mathcal M=\Psi^{\otimes n}$, always finite dimensional.
For an ensemble $\{p_x,\rho^x_{A^nB}\}$ and surviving subset $J$,
the output is $(\id_{A_J}\otimes\mathcal M)(\rho^x_{A_JB})$,
with entropy at least $\widehat S_{\mathcal M}(\rho^x_B)$.
Hence each full output has entropy at least
$nh(\lambda)+\widehat S_{\mathcal M}(\rho^x_B)$.
Subadditivity bounds the average output entropy by
\[
 nh(\lambda)+\lambda\sum_iS(\bar\rho_{A_i})
                         +S(\mathcal M(\bar\rho_B)).
\]
Subtract the conditional entropies and use roof convexity to obtain
\begin{align*}
\chi&\le\lambda\sum_iS(\bar\rho_{A_i})
       +S(\mathcal M(\bar\rho_B))-\widehat S_{\mathcal M}(\bar\rho_B)\\
    &\le n\lambda g(N)+\chi(\Psi^{\otimes n}).
\end{align*}
The constrained coding theorem and regularization \cite{holevo} yield
the desired upper bound. This direct subset argument avoids iteration
over infinite-dimensional partner channels.
Postprocessing gives the reset upper bound, and independent product
coding gives the lower bound
$C(\Rcal\otimes\Psi,N)\ge C(\Rcal,N)+C(\Psi)$.
Subtract and use Theorem~\ref{thm:main}.
\end{proof}
This is an asymptotic classical statement, not exact additivity at finite
energy. No quantum or private arbitrary-partner additivity is asserted.

\section{A sharp spectral criterion for absolute recovery}\label{sec:spectral}
The oscillator construction extends to a larger class of Hamiltonians,
but not to every Hamiltonian with finite entropy at finite energy.
In this section $H$ has a countably infinite eigenbasis, eigenvalues
$0=e_0\le e_1\le\cdots\to\infty$, finite multiplicities, and
$Z(\beta):=\Tr 2^{-\beta H}<\infty$ for at least one $\beta>0$.
Define
\[
 \beta_c=\inf\{\beta\ge0:Z(\beta)<\infty\},\qquad
 Z_c=\lim_{\beta\downarrow\beta_c}Z(\beta),\qquad
 F_H(E)=\sup_{\Tr H\rho\le E}S(\rho).
\]
Thus $\beta_c<\infty$ and $1<Z_c\le\infty$. The base of the
exponential fixes the units of $\beta$; all entropies remain in bits.
The input, error and security conventions are those above, with $H$
replacing the number operator. In particular, the two-way convention is
countably decomposable LOCC, with expected block-average input cost.
Fix $0<\lambda<1$ and any density operator $\tau$; neither its energy
nor its entropy is required to be finite.

\Needspace{18\baselineskip}
\begin{theorem}[Critical-partition dichotomy]\label{thm:spectral}
If $Z_c=\infty$, then for every fixed $\tau$,
\begin{equation}\label{eq:spectral-six}
 X(\Rcal,H,E)=r_X F_H(E)+o(1),\qquad E\to\infty,
\end{equation}
simultaneously for
\[
 (X,r_X)=
 (C,\lambda),\ (C_{\rm EA},2\lambda),\
 (Q,(2\lambda-1)_+),\ (P,(2\lambda-1)_+),\
 (Q_2,\lambda),\ (K_2,\lambda).
\]
Each benchmark is an upper bound. If $Z_c<\infty$, put
$\gamma_c=2^{-\beta_cH}/Z_c$. Then, for every fixed $\tau$,
\begin{align}
 2\lambda F_H(E)-C_{\rm EA}(\Rcal,H,E)
   &\longrightarrow\Delta_{\lambda,\tau}(H)>0,
       \label{eq:spectral-defect}\\
 \Delta_{\lambda,\tau}(H)
   &=\inf_{\rho:\Tr H\rho<\infty}
       \{2\lambda D(\rho\Vert\gamma_c)+e_{\lambda,\tau}(\rho)\},
       \label{eq:critical-variational}\\
 0<\Delta_{\lambda,\tau}(H)
   &\le e_{\lambda,\tau}(\gamma_c)\le h(\lambda).
       \label{eq:critical-bounds}
\end{align}
Here $e_{\lambda,\tau}$ is the finite binary-flag loss defined in
Lemma~\ref{lem:flag-continuity}. In particular, for each fixed reset
state, absolute entanglement-assisted recovery holds if and only if
$Z_c=\infty$. No finite-$Z_c$ formula for $Q,P,Q_2,K_2$ is asserted.
\end{theorem}

The critical Gibbs entropy description used below is established:
Shirokov \cite[Proposition 1a]{shirokoventropy} gives the variational
formula, critical slope, linear regime and possible nonattainment.
We use that description to distinguish two operational capacity regimes.
The general energy-capacity formulas of \cite{wq,dsw} with an
all-temperature Gibbs hypothesis cannot simply be applied here, since
$\beta_c$ may be positive. We give the erasure converse directly.

\subsection{Entropy escape and the erasure benchmark}
\begin{lemma}\label{lem:spectral-escape}
Under the stated spectral assumptions,
\begin{equation}\label{eq:critical-entropy}
 \frac{F_H(E)}{E}\longrightarrow\beta_c,\qquad
 F_H(E)-\beta_cE\nearrow\log_2 Z_c.
\end{equation}
If $Z_c=\infty$, there exist diagonal finite-rank states $\sigma_E$
and their support projections $P_E$ such that
\begin{equation}\label{eq:spectral-escape}
 \Tr H\sigma_E<E,\qquad
 S(\sigma_E)=F_H(E)-o(1),\qquad
 \Tr P_E\tau\longrightarrow0
\end{equation}
for every fixed trace-class density $\tau$.
\end{lemma}
\begin{proof}
For completeness, the relevant entropy argument can be expressed by
the Gibbs duality
\[
 \log_2 Z(\beta)=\sup_{E\ge0}\{F_H(E)-\beta E\}.
\]
The finite concave nondecreasing function $F_H$ has an asymptotic
slope. Gibbs' bound makes it at most $\beta_c$. A strictly smaller
slope would make the displayed supremum finite for some
$\beta<\beta_c$, a contradiction. All secant slopes dominate the
asymptotic slope, so $F_H(E)-\beta_cE$ is nondecreasing. Its
supremum is $\log_2 Z_c$, proving \eqref{eq:critical-entropy}.
This reasoning does not assume that the entropy supremum is attained.

Remove the first $m$ energy eigenvectors, and write $Z_{>m}$ and
$F_{>m}$ for the partition sum and entropy supremum on the remaining
space. For fixed $m$ and sufficiently large $E$, its Gibbs distribution
has mean $E$ at an inverse temperature $\beta_m(E)>\beta_c$.
Its mean tends to infinity as $\beta\downarrow\beta_c$: otherwise
integrating the bounded mean, the negative derivative of
$\log Z_{>m}$, would contradict the divergence of $Z_{>m}$.
Thus $\beta_m(E)\downarrow\beta_c$. Evaluating the full Gibbs bound
at that temperature gives
\[
 0\le F_H(E)-F_{>m}(E)
 \le\log_2\frac{Z(\beta_m(E))}{Z_{>m}(\beta_m(E))}
 \longrightarrow0.
\]
The removed partition sum is bounded and the retained sum diverges.
Choose $m=m(E)\to\infty$ slowly by a diagonal argument. Upper
truncation of each tail Gibbs distribution preserves its entropy in
the limit and lowers its mean energy. If necessary, first decrease
the chosen mean by a positive amount tending sufficiently quickly
to zero; continuity of the finite concave entropy supremum supplies
strict energy slack with $o(1)$ entropy loss. The resulting finite
spectral window lies above its first $m(E)$ vectors. Consequently
$\Tr P_E\tau\le\Tr\mathbf1_{\{j>m(E)\}}\tau\to0$ for every fixed
$\tau$. There is no uniform rate over all density operators.
\end{proof}

\begin{lemma}[Erasure with a single finite partition sum]\label{lem:general-erasure}
For the six operational conventions of Theorem~\ref{thm:spectral},
\begin{equation}\label{eq:general-erasure}
 X(\Ecal,H,E)=r_XF_H(E).
\end{equation}
\end{lemma}
\begin{proof}
For $E>0$, finite-support entropy approximants with strict cost slack, followed
by finite-dimensional erasure coding, give achievability. The order
is: fix the finite approximation, take block length to infinity,
then improve the approximation. For two-way transmission the
receiver announces the surviving copies and distills their pure
entanglement. No entropy-maximizing infinite-dimensional code state
is required. At $E=0$, use the finite ground eigenspace directly.

For the converse, any admissible $n$-use average input has
$S(\rho_{A^n})\le\sum_iS(\rho_{A_i})\le nF_H(E)<\infty$.
The erasure outputs and complementary outputs also have finite
entropy, bounded by this sum plus $nh(\lambda)$.
Conditioning on the independent erasure pattern bounds the output
Holevo quantity by $\lambda\sum_i S(\rho_{A_i})$.
Fano's inequality on the finite message alphabet proves the $C$
converse. For assistance, $I(\rho,\Ecal)=2\lambda S(\rho)$;
mutual-information subadditivity gives the corresponding $n$-use
bound. The general constrained entanglement-assisted formula
\cite[Theorem 1]{hs} applies to the present positive observable.

For $\lambda\le1/2$ erasure is antidegradable, so $Q=P=0$.
For $\lambda>1/2$ let $B^n$ and $W^n$ be its output and full
environment. Product degradability and data processing of total
correlation give
\begin{align*}
 I_c(\rho,\Ecal^{\otimes n})
 &=\sum_i I_c(\rho_{A_i},\Ecal)
       -\operatorname{TC}(B^n)+\operatorname{TC}(W^n)\\
 &\le(2\lambda-1)\sum_i S(\rho_{A_i}).
\end{align*}
Here $\operatorname{TC}(B^n)=\sum_iS(B_i)-S(B^n)$, and similarly
for $W^n$. All these entropy expressions are finite. If a quantum encoder has
an additional purifying system $F$, then
$S(F)\le S(R)+S(A^n)<\infty$ for its finite logical reference $R$.
The complementary channel is antidegradable, so $FW^n$ has a
two-extension on $W^n$. Weak monotonicity gives $S(F|W^n)\ge0$,
and purity of the full dilation gives
$I(R\rangle B^n)=I_c(\rho,\Ecal^{\otimes n})-S(F|W^n)$.
Coherent-information monotonicity under decoding and continuity on
the finite logical reference prove the quantum converse.
For any finite private ensemble, its private information is
$I_c(\bar\rho)-\sum_xp_x I_c(\rho_x)$.
Degradability makes each subtracted coherent information nonnegative;
classical Fano and secrecy continuity on the finite message register
give the same private converse. None of these arguments uses uniform
entropy continuity on the whole energy ball.

For two-way protocols consider one input system $A$ held by Alice,
with her remaining memory $L$ and Bob's memory $B$. Transferring $A$
to Bob can increase squashed entanglement by at most $S(A)$.
Indeed the extended conditional-mutual-information chain rule
\cite[equation (14)]{shirokovcontinuity} gives, for any extension $W$,
\[
 I(L;AB|W)\le I(LA;B|W)+2S(A).
\]
Choose an extension within any positive tolerance of the pre-use
squashed entanglement. Starting from zero, this quantity remains
finite by induction, and $S(A)$ is finite by the input cost, so
there is no subtraction of infinite quantities. Failure is a local
discard and cannot increase squashed entanglement. Store the
independent success flag in the extension; the expected increment
is at most $\lambda S(A)$. Average over public histories and apply
entropy concavity and Jensen's inequality to obtain
\[
 E_{\rm sq}(\text{final state})\le\lambda nF_H(E).
\]
We use squashed entanglement in infinite dimensions
\cite[Proposition 2B]{shirokovsquashed} and the approximate-private-state bound
\cite[Theorem 2]{wildeprivate}. Retain the local purifiers
as shields when representing the key protocol by an approximate private
state; Eve includes the public transcript. The bound is applied to
that state, before measuring its key registers. Its error terms
depend on the finite key size, not on shield dimension.
With vanishing error it bounds the secret-key rate, and therefore
the entanglement rate, by $\lambda F_H(E)$.
The operations retain the stated countably decomposable convention.
\end{proof}

\begin{proof}[Proof of \eqref{eq:spectral-six}]
Decode the finite spectral window of Lemma~\ref{lem:spectral-escape}
by projection and an orthogonal failure flag. The accepted reset
weight is $\delta_E=(1-\lambda)\Tr P_E\tau\to0$ and its binary
ambiguity cost is $b_E=(\lambda+\delta_E)
h(\lambda/(\lambda+\delta_E))\to0$.
The finite-channel results proved above give respectively the lower
bounds
\[
 \lambda S(\sigma_E)-b_E,\quad
 2\lambda S(\sigma_E)-b_E,\quad
 [(2\lambda-1)S(\sigma_E)-b_E]_+,
\]
for $C,C_{\rm EA},Q/P$, and
$[\lambda S(\sigma_E)-b_E-\delta_E\log_2 5]_+$ for $Q_2/K_2$.
The paid seed uses copies of the same strictly admissible input.
Reset is postprocessing of erasure, so Lemma~\ref{lem:general-erasure}
supplies all upper bounds. Substitution of
$S(\sigma_E)=F_H(E)-o(1)$ proves the result. For $\lambda\le1/2$,
the quantum and private capacities are identically zero.
\end{proof}

\subsection{A finite flag and a positive variational defect}
The necessity statement is stronger than a counterexample using a
vacuum reset: every fixed reset state has a positive limiting
entanglement-assisted deficit when $Z_c$ is finite.

\begin{lemma}\label{lem:flag-continuity}
For any density operators $\rho,\tau$, choose a purification
$\psi_\rho$ with $B$ marginal $\rho$, and let
\[
 \Omega_{JRB}=
 \lambda|0\rangle\langle0|_J\otimes\psi_\rho
 +(1-\lambda)|1\rangle\langle1|_J\otimes\rho_R\otimes\tau_B.
\]
The quantity $e_{\lambda,\tau}(\rho):=I(J;R|B)_\Omega$ is
well defined, continuous in trace norm in $\rho$ and $\tau$, and
lies in $[0,h(\lambda)]$. If $S(\rho)<\infty$, then
\begin{equation}\label{eq:flag-loss-exact}
 I(\rho,\Rcal)=2\lambda S(\rho)-e_{\lambda,\tau}(\rho).
\end{equation}
For every faithful mixed state $\gamma$,
$e_{\lambda,\tau}(\gamma)>0$, even if $S(\gamma)=\infty$.
\end{lemma}
\begin{proof}
Equivalently, $e$ is the difference of the Holevo quantities of
the two binary ensembles
$\{\lambda,\psi_\rho;1-\lambda,\rho_R\otimes\tau\}$ and
$\{\lambda,\rho;1-\lambda,\tau\}$.
Define those quantities through relative entropy. Both are finite
and at most $h(\lambda)$, regardless of the state entropies.
Data processing gives the asserted bounds. Keeping $J$ in the
channel output gives $I(R;BJ)=2\lambda S(\rho)$; the chain rule
then proves \eqref{eq:flag-loss-exact} when $S(\rho)$ is finite.
Canonical purifications vary continuously in trace norm, by the
square-root continuity of positive trace-class operators. Conditional
mutual information is uniformly continuous when one unconditioned
system is the two-dimensional $J$
\cite[Corollary 1, equation (25)]{shirokovcontinuity}.
This proves continuity of $e$ without continuity of $S(\rho)$.

Suppose $e_{\lambda,\tau}(\gamma)=0$. The binary Holevo data
processing inequality under $\Tr_R$ is then an equality.
Each component has finite relative entropy to its ensemble average,
because that average dominates it by its positive prior weight.
Equality in relative-entropy data processing therefore supplies one
normal recovery channel $T:B\to RB$ with
\[
 T(\gamma)=\psi_\gamma,\qquad
 T(\tau)=\gamma_R\otimes\tau.
\]
Here we use the infinite-dimensional sufficiency theorem
\cite[Section 2.3, Theorem 2]{jencova}, restricted to the support
of the ensemble average; no bounded density inverse is assumed.
A channel mapping a faithful state to a pure state is constant:
if $Q=I-\psi_\gamma$, then
$\Tr\gamma T^*(Q)=0$ and positivity and faithfulness imply
$T^*(Q)=0$. Every output is consequently supported on that same
one-dimensional subspace. Thus $T(\tau)=\psi_\gamma$, contrary
to the mixed product state $\gamma_R\otimes\tau$.
\end{proof}

\begin{proof}[Proof of \eqref{eq:spectral-defect}--\eqref{eq:critical-bounds}]
If $Z_c<\infty$, then $\beta_c>0$ and $\gamma_c$ is faithful.
For every finite-energy input, entropy is finite and the Gibbs
identity gives
\[
 D(\rho\Vert\gamma_c)
   =\beta_c\Tr H\rho+\log_2 Z_c-S(\rho)<\infty.
\]
The constrained entanglement-assisted formula \cite[Theorem 1]{hs}
and Lemma~\ref{lem:flag-continuity} imply
\[
 2\lambda F_H(E)-h(\lambda)
 \le C_{\rm EA}(\Rcal,H,E)\le2\lambda F_H(E).
\]
Write this finite concave nondecreasing capacity as $C_E$. By
\eqref{eq:critical-entropy}, $C_E/E\to2\lambda\beta_c$.
Concavity implies that $C_E-2\lambda\beta_c E$ is nondecreasing,
so its limit equals its supremum. Optimizing the energy allowance
down to the mean cost of each input yields
\begin{align*}
 \sup_E(C_E-2\lambda\beta_c E)
 &=\sup_{\Tr H\rho<\infty}
       [I(\rho,\Rcal)-2\lambda\beta_c\Tr H\rho]\\
 &=2\lambda\log_2 Z_c-\Delta_{\lambda,\tau}(H).
\end{align*}
This identity uses suprema, not existence of an optimizing state.
Combining it with \eqref{eq:critical-entropy} gives the limit.

If the infimum in \eqref{eq:critical-variational} were zero,
both its nonnegative summands would vanish along a sequence.
Pinsker's inequality would give $\rho_n\to\gamma_c$ in trace
norm. Continuity would then give
$e_{\lambda,\tau}(\rho_n)\to e_{\lambda,\tau}(\gamma_c)>0$,
a contradiction. For the upper bound take normalized finite
energy truncations $\gamma_{c,m}$. Their relative entropy is
$-\log_2\Tr(P_m\gamma_c)\to0$, and their flag loss converges
to that of $\gamma_c$. This remains valid when the critical
state has infinite mean energy or infinite entropy. Neither
attainment of the variational infimum nor a positive lower bound
uniform over all reset states is claimed.
\end{proof}

\subsection{An explicit permanent loss at arbitrarily large energy}
For a unique ground state and the reset $\tau=|0\rangle\langle0|$,
the defect is a scalar optimization. Put $q=1-\lambda$ and define
\begin{align}
 a&=\lambda+qx,\qquad
 \mu_\pm=\tfrac12\bigl(a\pm\sqrt{a^2-4\lambda qx(1-x)}\bigr),\nonumber\\
 f_\lambda(x)&=h(x)+h(\lambda(1-x))
       -H(\mu_+,\mu_-,q(1-x)),\label{eq:vacuum-scalar}\\
 e_\lambda(x)&=2\lambda h(x)-f_\lambda(x).\nonumber
\end{align}
For a diagonal input $x|0\rangle\langle0|+(1-x)\sigma$,
$I(\rho,\Rcal)=2\lambda(1-x)S(\sigma)+f_\lambda(x)$.
These eigenvalues and this decomposition are the established
vacuum-reset calculation of Catalano et al.
\cite[Section VI.B and Appendix C, equations (C1)--(C11)]{catalano}.
For $E>0$, applying it to the excited-space entropy supremum $F_{H,+}$ gives
\begin{equation}\label{eq:vacuum-general-H}
 C_{\rm EA}(\Rcal,H,E)=
 \sup_{0\le x<1}
 \{2\lambda(1-x)F_{H,+}(E/(1-x))+f_\lambda(x)\},
\end{equation}
where $x$ must meet the minimum excited-energy constraint; the
actual $x=1$ input gives $C_{\rm EA}(\Rcal,H,0)=0$ at $E=0$.
To justify diagonal restriction without uniform energy-ball entropy
continuity, average successively over phases on the first finitely
many energy vectors. Covariance and concavity do not decrease mutual
information. The pinched inputs converge to the diagonal input,
and their entropies increase to its finite diagonal entropy, by
lower semicontinuity and the entropy upper bound on that diagonal.
The same pinching argument applies to output and complement; the
latter retains the displayed two-dimensional block and the remaining
eigenvalues $q(1-x)\operatorname{spec}(\sigma)$. Their limiting
entropies are finite. Passing to the limit preserves the inequality.

If $Z_c<\infty$, let $Z_+=Z_c-1$ and $p_c=1/Z_c$.
The excited-space version of \eqref{eq:critical-entropy}, substituted
into \eqref{eq:vacuum-general-H}, gives
\begin{equation}\label{eq:vacuum-defect}
 \Delta_{\lambda,|0\rangle}(H)=
 \min_{0\le x\le1}
       \{2\lambda D_{\rm Bern}(x\Vert p_c)+e_\lambda(x)\}.
\end{equation}
The remainder is uniform because $E/(1-x)\ge E$; the $x\to1$
endpoint is approached through admissible $x<1$, and is not a
claim that the vacuum carries energy. Strict positivity can also
be checked directly: dephase and coarsen the reference to its
ground/excited label in the flag-loss expression. This gives
\[
 e_\lambda(x)\ge
 (q+\lambda x)h\!\left(\frac q{q+\lambda x}\right)-xh(\lambda)>0
 \quad(0<x<1).
\]
The last expression is a classical conditional mutual information:
at the vacuum output an excited reference forces failure, whereas
a ground reference allows either branch. Compactness and the
strictly interior $p_c$ prove positivity in \eqref{eq:vacuum-defect}.

\begin{proposition}[An exact constant-loss plateau]\label{prop:spectral-example}
Take one ground vector and, for each integer $n\ge1$, a shell of
energy $n^2+n$ with multiplicity $2^{n^2}$. Then
$\beta_c=1$, $Z_c=2$, and, for the ground-state reset and every $E\ge8$,
\begin{align}
 C_{\rm EA}(\Rcal,H,E)&=2\lambda E+\max_{0\le x\le1}f_\lambda(x),
       \label{eq:plateau-capacity}\\
 2\lambda F_H(E)-C_{\rm EA}(\Rcal,H,E)
       &=2\lambda-\max_x f_\lambda(x)>0.\label{eq:plateau-gap}
\end{align}
The scalar maximum is unique and interior. At $\lambda=1/2$ the
gap is numerically $0.409719932065909\ldots$ bits per use.
The capacity equality is a supremum over operational codes; an
entropy-maximizing finite-energy input need not exist.
\end{proposition}
\begin{proof}
Here
$Z(\beta)=1+\sum_{n\ge1}2^{(1-\beta)n^2-\beta n}$,
which establishes the critical parameters. The normalized excited
critical Gibbs law has shell probabilities $2^{-n}$, mean energy
$\sum(n^2+n)2^{-n}=8$ and entropy $8$. The full critical state
has ground probability $1/2$, mean energy $4$ and entropy $5$.
The Gibbs entropy suprema \cite[Proposition 1a]{shirokoventropy}
are therefore
\[
 F_H(E)=E+1\quad(E\ge4),\qquad
 F_{H,+}(E)=E\quad(E\ge8).
\]
One can see the nonattained linear part explicitly: truncate the
critical Gibbs law, then mix in a uniform state on a much higher
shell $m$. Its weight is of order the extra energy divided by
$m^2+m$, while its entropy per added energy tends to
$m^2/(m^2+m)\to1$. Let the upper shell tend to infinity and
then improve the finite base truncation. The displayed upper
bounds are approached. Above the critical mean, equality in the
Gibbs bound would force the critical Gibbs state itself, which
has the wrong energy; thus these are suprema rather than maxima.
Substitution into \eqref{eq:vacuum-general-H} proves
\eqref{eq:plateau-capacity}, and \eqref{eq:vacuum-defect} proves
the positive gap.

Finally $f_\lambda(x)$ is the mutual information of a two-level
diagonal input. For any nontrivial mixture of two different such
inputs, its concavity gap is
$\chi_{\rm input}+\chi_{\rm output}-\chi_{\rm environment}
\ge\chi_{\rm output}>0$, by data processing and $\lambda>0$.
Hence it is strictly concave. It vanishes at both endpoints and
is positive in the interior, as is already seen by measuring
the reference in the binary basis. The reported numerical root
is a diagnostic evaluation of this unique maximum, not a proof
by sampling.
\end{proof}

For contrast, a shell of energy $n$ and multiplicity $2^n$ for
each $n\ge0$ has $\beta_c=1$ but $Z_c=\infty$, and
$F_H(E)=E+g(E)$. Removing the first $m$ shells leaves the entropy
supremum $E+g(E-m)$ for $E\ge m$; its deficit tends to zero.
Thus a positive critical inverse temperature is not itself the
obstruction. Normalizability at that temperature is decisive.
These deliberately constructed spectra establish the mathematical
boundary of the capacity law; no experimental Hamiltonian or
thermodynamic phase transition is claimed.


\section{Consequences and quantitative checks}
\begin{corollary}
For every $N>0$, $Q(\Rcal,N)>0$ and $P(\Rcal,N)>0$
if and only if $\lambda>1/2$.
\end{corollary}
\begin{proof}
Only positivity needs proof. If $\lambda>1/2$, the main theorem gives
a finite $E_0$ and a positive quantum rate. Use those codes on a fixed
fraction $0<\theta<\min\{1,N/E_0\}$ of block positions and send vacuum
on the rest. The rate remains positive and mean energy is below $N$.
A quantum code yields a private code.
\end{proof}
This uses block-average energy and does not assert a peak-energy result.

\begin{corollary}
For every oscillator density operator $\tau$, even without a finite mean
energy assumption, $\Rcal$ is not entanglement breaking if $\lambda>0$.
\end{corollary}
\begin{proof}
Fix $d=2$ and let $K\to\infty$. Trace class implies $\epsilon\to0$.
Apply the decoded channel to half of a normalized Bell pair.
In the transmitted block, the expectation of the partial transpose on
the normalized antisymmetric vector is $-\lambda/2+\delta/4$,
which is negative for large $K$.
Thus the decoded channel has an entangled Choi state, impossible if
the original channel were entanglement breaking.
\end{proof}
The vacuum versions overlap \cite{catalano}; these structural consequences
are not presented as the main contribution.

\paragraph{Numerical checks.}
The original accompanying script checks 200 random finite-matrix cases, including
reset coherences and entangled references.
The largest discrepancies in the distance identity and erasure coherent
information were $4.45\,10^{-16}$ and $2.34\,10^{-15}$.
No tested finite-energy, continuity or classical lower bound failed.
A separate new-seed script checks 500 finite-matrix cases for the
hidden-branch, forward/reverse hashing and pinching inequalities,
including pure reset states, zero overlap and survival endpoints.
The largest negative residual in those lower-bound inequalities has
magnitude below $1.2\,10^{-15}$, consistent with floating-point error.
The partial-breeding diagnostics add 120 cases checking parity Kraus
operators, the fixed-matrix entropy bound and survivor hashing.
A further 240 cases test the full phase average and general-noise bounds,
including 120 separable states with matching marginals. Their algebraic
errors are below $4.1\,10^{-15}$ and no lower-bound violation exceeds
rounding error. A further 1000 cases check the dominated-component
seed algebra without matching marginals, including singular supports;
480 cases check mixed maximally correlated components and survivor
entropies. Another 264 cases check invariant-state block entropies and
finite phase averages, including explicit control families. These
finite-matrix checks are not proofs of asymptotic coding,
security, capacity, or exhaustive counterexample searches.
The minimax supplement adds 65 exact rational checks and 45 numerical
parameter cases, including an independent nonidentity Weyl construction,
correlated-block spectra and the interval-width identity; its largest
algebraic residual is below $2\,10^{-14}$. It does not test coding
or privacy by simulation.

Table~\ref{tab:bounds} gives conservative analytic gap bounds from
\eqref{eq:choose} for $\lambda=0.7$ and only the assumption $M\le10$.
They bound absolute losses from the erasure benchmarks, not measured
rates or optimized capacity estimates.
\begin{table}[ht]
\centering
\small\begin{tabular}{rrrrr}
\toprule
$N$&Classical&Quantum/private&EA classical&Two-way $Q_2,K_2$\\
\midrule
$10^{2}$&0.988103&0.942327&1.094915&1.684682\\
$10^{3}$&0.448692&0.434617&0.481535&0.666373\\
$10^{4}$&0.190675&0.186325&0.200824&0.260333\\
$10^{5}$&0.075710&0.074336&0.078917&0.097685\\
$10^{6}$&0.028947&0.028514&0.029957&0.035913\\
$10^{8}$&0.003891&0.003848&0.003992&0.004587\\
\bottomrule
\end{tabular}
\caption{Analytic upper bounds on gaps, rounded upward to six decimals.
A loose bound can exceed the benchmark itself, in which case a trivial
zero-rate lower bound is stronger.}\label{tab:bounds}
\end{table}

\paragraph{Limitations.}
The reset state is fixed; no uniform convergence over all density
operators is claimed. The pure-component irreversibility theorem requires separable noise
for its cost upper bound, but no matching marginal. A mixed maximally
correlated component can have intrinsic cost irreversibility. The
CLDUI/LDOI upper bounds apply to those invariant classes; only the
distillation lower bound applies before twirling an arbitrary state.
The argument does not cover reset states changing arbitrarily with $N$,
general infinite-dimensional input-dependent noise, partial-loss fading
components or Gaussian thermal attenuation. The separate finite-dimensional
arbitrary-noise result in Section~\ref{sec:minimax} concerns fixed memoryless
channels and codes that can depend on them. It does not provide adversarial
noise protection or evaluate the unknown depolarizing capacity functions.
The oscillator result determines neither general finite-$N$ capacities nor optimal error
rates or second-order coding costs. The growing code dimension is an
asymptotic construction, not an efficient experimental receiver.
Exact finite-energy classical additivity for arbitrary reset states
would require a separate argument.

\paragraph{Spectral checks.}
The spectral checks independently evaluate the vacuum joint-state
spectrum on 200 diagonal inputs, test phase pinching on 200 complex
inputs, and check 2,807 scalar positivity cases. Twenty-five parameter
choices compare the two defect expressions; five plateau maxima are
evaluated with 80-digit arithmetic. Forty-two finite high-shell mixtures
check the nonattained entropy construction and fifteen cases check
the contrasting divergent-partition example. These are finite numerical
diagnostics, not proofs of the spectral dichotomy, capacity converses,
infinite-dimensional recovery or priority. The analytic arguments above
supply those mathematical claims, subject to their cited dependencies.

\paragraph{Provenance and reproducibility.}
This working paper was developed with AI assistance in exploration,
drafting and computational checking. A model review is not human peer
review or formal proof verification. The supplement contains the
verification scripts and numerical records; external source PDFs and
private review records are excluded. Erasure coding, the energy-constrained
coding theorems, hashing and pure-state dilution are cited dependencies,
distinct from the hidden-branch estimate, stopped-classification bound
and their numerical checks. No model score is evidence of a theorem.

\begin{thebibliography}{99}
\bibitem{bds} C.~H. Bennett, D.~P. DiVincenzo and J.~A. Smolin,
Capacities of quantum erasure channels,
\emph{Phys.\ Rev.\ Lett.} \textbf{78}, 3217--3220 (1997).
\href{https://arxiv.org/abs/quant-ph/9701015}{arXiv:quant-ph/9701015}.
\bibitem{winter} A. Winter, Tight uniform continuity bounds for quantum
entropies: conditional entropy, relative entropy distance and energy
constraints, \emph{Commun.\ Math.\ Phys.} \textbf{347}, 291--313 (2016).
\href{https://arxiv.org/abs/1507.07775}{arXiv:1507.07775}.
\bibitem{wq} M.~M. Wilde and H. Qi, Energy-constrained private and quantum
capacities of quantum channels, \emph{IEEE Trans.\ Inf.\ Theory}
\textbf{64}, 7802--7827 (2018).
\href{https://arxiv.org/abs/1609.01997}{arXiv:1609.01997}.
\bibitem{holevo} A.~S. Holevo, Entanglement-Assisted Capacities of Constrained
Quantum Channels, \emph{Theory Probab.\ Appl.} \textbf{48}(2), 243--255 (2004),
\href{https://doi.org/10.1137/S0040585X97980415}{doi:10.1137/S0040585X97980415}.
\href{https://arxiv.org/abs/quant-ph/0211170}{arXiv:quant-ph/0211170}.
\bibitem{hs} A.~S. Holevo and M.~E. Shirokov, On classical capacities of
infinite-dimensional quantum channels,
\emph{Problems of Information Transmission} \textbf{49}, 15--31 (2013).
\href{https://doi.org/10.1134/S003294601301002X}{doi:10.1134/S003294601301002X}.
Preprint title: On the entanglement-assisted classical capacity of
infinite-dimensional quantum channels,
\href{https://arxiv.org/abs/1210.6926}{arXiv:1210.6926}.
\bibitem{zoj} C. Zhong, C. Oh and L. Jiang, Information transmission with
continuous variable quantum erasure channels,
\href{https://arxiv.org/abs/2205.09711}{arXiv:2205.09711} (2022).
\bibitem{catalano} G. Catalano, M. Fanizza, F.~A. Mele, G. De Palma and
V. Giovannetti, Convex combinations of bosonic pure-loss channels,
\href{https://arxiv.org/abs/2604.26874v2}{arXiv:2604.26874v2}
(30 September 2026).
\bibitem{amosov} G.~G. Amosov and S. Mancini, The decreasing property of
relative entropy and the strong superadditivity of quantum channels,
\href{https://arxiv.org/abs/0803.0452}{arXiv:0803.0452} (2008).
\bibitem{dsw} N. Davis, M.~E. Shirokov and M.~M. Wilde,
Energy-constrained two-way assisted private and quantum capacities of
quantum channels, \emph{Phys.\ Rev.\ A} \textbf{97}, 062310 (2018).
\href{https://arxiv.org/abs/1801.08102v3}{arXiv:1801.08102v3}.
\bibitem{dw} I. Devetak and A. Winter, Distillation of secret key and
entanglement from quantum states,
\href{https://arxiv.org/abs/quant-ph/0306078}{arXiv:quant-ph/0306078} (2003).
\bibitem{bbps} C.~H. Bennett, H.~J. Bernstein, S. Popescu and B. Schumacher,
Concentrating partial entanglement by local operations,
\emph{Phys.\ Rev.\ A} \textbf{53}, 2046--2052 (1996).
\href{https://arxiv.org/abs/quant-ph/9511030}{arXiv:quant-ph/9511030}.
\bibitem{hostens} E. Hostens, J. Dehaene and B. De Moor,
Asymptotic adaptive bipartite entanglement-distillation protocol,
\emph{Phys.\ Rev.\ A} \textbf{73}, 062337 (2006).
\href{https://arxiv.org/abs/quant-ph/0602205v2}{arXiv:quant-ph/0602205v2}.
\bibitem{vv} K.~G.~H. Vollbrecht and F. Verstraete,
Interpolation of recurrence and hashing entanglement distillation protocols,
\emph{Phys.\ Rev.\ A} \textbf{71}, 062325 (2005).
\href{https://arxiv.org/abs/quant-ph/0404111v2}{arXiv:quant-ph/0404111v2}.
\bibitem{vw} K.~G.~H. Vollbrecht and M.~M. Wolf,
Efficient distillation beyond qubits,
\emph{Phys.\ Rev.\ A} \textbf{67}, 012303 (2003).
\href{https://arxiv.org/abs/quant-ph/0208152}{arXiv:quant-ph/0208152}.
\bibitem{hh} M. Horodecki and P. Horodecki,
Reduction criterion of separability and limits for a class of distillation
protocols, \emph{Phys.\ Rev.\ A} \textbf{59}, 4206--4216 (1999).
\href{https://arxiv.org/abs/quant-ph/9708015v3}{arXiv:quant-ph/9708015v3}.
\bibitem{gew} K. Goodenough, D. Elkouss and S. Wehner,
Assessing the performance of quantum repeaters for all phase-insensitive
Gaussian bosonic channels, \emph{New J.\ Phys.} \textbf{18}, 063005 (2016).
\href{https://arxiv.org/abs/1511.08710v2}{arXiv:1511.08710v2}.
\bibitem{maneva} E.~N. Maneva and J.~A. Smolin,
Improved two-party and multi-party purification protocols,
\href{https://arxiv.org/abs/quant-ph/0003099}{arXiv:quant-ph/0003099} (2000).
\bibitem{leung} A.~W. Leung and P.~W. Shor,
Adaptive Entanglement Purification Protocols with Two-way Classical
Communication,
\href{https://arxiv.org/abs/quant-ph/0702156v3}{arXiv:quant-ph/0702156v3} (2007).
\bibitem{hhho} K. Horodecki, M. Horodecki, P. Horodecki and J. Oppenheim,
General paradigm for distilling classical key from quantum states,
\href{https://arxiv.org/abs/quant-ph/0506189}{arXiv:quant-ph/0506189} (2005).
\bibitem{vwentropy} K.~G.~H. Vollbrecht and M.~M. Wolf,
Conditional entropies and their relation to entanglement criteria,
\href{https://arxiv.org/abs/quant-ph/0202058}{arXiv:quant-ph/0202058} (2002).
\bibitem{jm} N. Johnston and O. MacLean,
Pairwise completely positive matrices and conjugate local diagonal
unitary invariant quantum states,
\emph{Electronic Journal of Linear Algebra} \textbf{35}, 156--180 (2019).
\href{https://arxiv.org/abs/1807.06897}{arXiv:1807.06897}.
\bibitem{sn} S. Singh and I. Nechita,
Diagonal unitary and orthogonal symmetries in quantum theory,
\emph{Quantum} \textbf{5}, 519 (2021).
\href{https://doi.org/10.22331/q-2021-08-09-519}{doi:10.22331/q-2021-08-09-519}.
\bibitem{bl} J. Britz and M. Laurent,
Semidefinite hierarchies for diagonal unitary invariant bipartite
quantum states,
\href{https://arxiv.org/abs/2512.06551v2}{arXiv:2512.06551v2} (2026).

\bibitem{bdsw} C.~H. Bennett, D.~P. DiVincenzo, J.~A. Smolin and
W.~K. Wootters, Mixed-state entanglement and quantum error correction,
\emph{Phys.\ Rev.\ A} \textbf{54}, 3824--3851 (1996).
\href{https://arxiv.org/abs/quant-ph/9604024}{arXiv:quant-ph/9604024}.
\bibitem{king} C. King, The capacity of the quantum depolarizing channel,
\emph{IEEE Trans.\ Inf.\ Theory} \textbf{49}, 221--229 (2003).
\href{https://arxiv.org/abs/quant-ph/0204172}{arXiv:quant-ph/0204172}.
\bibitem{bsst} C.~H. Bennett, P.~W. Shor, J.~A. Smolin and
A.~V. Thapliyal, Entanglement-Assisted Classical Capacity of Noisy
Quantum Channels, \emph{Phys.\ Rev.\ Lett.} \textbf{83}, 3081--3084 (1999).
\href{https://arxiv.org/abs/quant-ph/9904023}{arXiv:quant-ph/9904023}.
\bibitem{plob} S. Pirandola, R. Laurenza, C. Ottaviani and L. Banchi,
Fundamental limits of repeaterless quantum communications,
\emph{Nature Communications} \textbf{8}, 15043 (2017).
\href{https://arxiv.org/abs/1510.08863}{arXiv:1510.08863}.
\bibitem{shirokoventropy} M.~E. Shirokov,
Entropic characteristics of subset of states,
\href{https://arxiv.org/abs/quant-ph/0510073v2}{arXiv:quant-ph/0510073v2}
(2005), Proposition 1a.
\bibitem{shirokovcontinuity} M.~E. Shirokov,
Tight continuity bounds for the quantum conditional mutual information,
for the Holevo quantity and for capacities of quantum channels,
\href{https://arxiv.org/abs/1512.09047v7}{arXiv:1512.09047v7} (2017).
\bibitem{jencova} A. Jen\v{c}ov\'a,
Recoverability of quantum channels via hypothesis testing,
\emph{Lett. Math. Phys.} \textbf{114}, 31 (2024),
\href{https://doi.org/10.1007/s11005-024-01775-2}{doi:10.1007/s11005-024-01775-2}.
\href{https://arxiv.org/abs/2303.11707}{arXiv:2303.11707}.
\bibitem{shirokovsquashed} M.~E. Shirokov,
Squashed entanglement in infinite dimensions,
\href{https://arxiv.org/abs/1507.08964v5}{arXiv:1507.08964v5} (2015),
Proposition 2B.
\bibitem{wildeprivate} M.~M. Wilde,
Squashed entanglement and approximate private states,
\href{https://arxiv.org/abs/1606.08028v3}{arXiv:1606.08028v3} (2018).

\end{thebibliography}
\end{document}
